\Needspace{5\baselineskip}
\SourceLeaf{422}%
\begin{center}6.\end{center}
\noindent\hspace*{1.2em}If we regard $x$ and $y$ as functions of $s$, the formulas take an even more elegant form. Set
\[
\begin{aligned}
\frac{dx}{ds}&=x',&\qquad \frac{dx'}{ds}&=x'',\\
\frac{dy}{ds}&=y',&\frac{dy'}{ds}&=y''.
\end{aligned}
\]
Then
\[
\begin{aligned}
x'&=\cos\varphi,&\qquad y'&=\sin\varphi,\\
x''&=-\frac{\sin\varphi}{r},&y''&=\frac{\cos\varphi}{r},
\end{aligned}
\]
or
\[
y'=-rx'',\qquad x'=ry''
\]
or equivalently
\[
1=r(x'y''-y'x''),
\]
so that
\[
x'y''-y'x''
\]
represents the curvature and
\[
\frac{1}{x'y''-y'x''}
\]
the radius of curvature.

\Needspace{5\baselineskip}
\begin{center}7.\end{center}
\noindent\hspace*{1.2em}We now proceed to consider curved surfaces. To visualize the directions of straight lines in three-dimensional space, we imagine a sphere of radius $1$ drawn about an arbitrary center. A point on the sphere then represents the direction of all straight lines parallel to the radius at whose endpoint the point lies. Just as the positions of all points in space are determined by their perpendicular distances $x,y,z$ from three mutually perpendicular planes, so the directions of the three principal axes perpendicular to those principal planes are to be marked on the surface of the auxiliary sphere\space%
